%% ==========================================================================
%% ADDENDUM to Part G (Certify v0.7.0, DOI 10.5281/zenodo.21890619),
%% "An Erdos--Rado oriented Ramsey number determined: k(3,4)=r(I_3,L_4)=21".
%% DRAFT -- August 11, 2026. Author byline: Daniel Kirtchakov (05oz,
%% halfounce.io; no institutional affiliation; ORCID 0009-0009-5213-4098).
%%
%% Answers Question 8.1 of Part G in the negative, at exactly the strength
%% independently certified here: at least THIRTEEN pairwise non-isomorphic
%% {I_3,TT_4}-free oriented graphs on 20 vertices exist, all rigid.
%% Novelty swept clean on the drafting date (2026-08-11): the value
%% k(3,4)=21 has no prior appearance beyond Part G (erdosproblems.com/112
%% OPEN; IRW the only paper on r(I_m,L_n)), so a fortiori its extremal
%% multiplicity is new.  Framing per PROTOCOL section 11: open with the
%% question, cite computation in one paragraph, close with open questions.
%% ==========================================================================
\documentclass[11pt]{amsart}

\usepackage[margin=1in]{geometry}
\usepackage{amsmath,amssymb}
\usepackage{booktabs}
\usepackage{array}
\usepackage[colorlinks=true,allcolors=blue]{hyperref}

\emergencystretch=3em
\tolerance=1200
\hbadness=4000

\newcommand{\TT}[1]{\mathrm{TT}_{#1}}
\newcommand{\I}[1]{I_{#1}}
\newcommand{\QR}{\mathrm{QR}_7}
\newcommand{\Aut}{\mathrm{Aut}}
\DeclareMathOperator{\Cay}{Cay}

\theoremstyle{plain}
\newtheorem{theorem}{Theorem}[section]
\newtheorem{lemma}[theorem]{Lemma}
\newtheorem{proposition}[theorem]{Proposition}
\newtheorem{corollary}[theorem]{Corollary}
\theoremstyle{definition}
\newtheorem{definition}[theorem]{Definition}
\theoremstyle{remark}
\newtheorem{remark}[theorem]{Remark}
\newtheorem{question}[theorem]{Question}

\begin{document}

\title[The $k(3,4)=21$ extremal graph is not unique]{The extremal graph
for $k(3,4)=r(I_3,L_4)=21$ is not unique: thirteen rigid witnesses, and a
forced Paley tournament}

\author{Daniel Kirtchakov}
\address{Independent researcher (\texttt{05oz}); no institutional
affiliation. ORCID:
\href{https://orcid.org/0009-0009-5213-4098}{0009-0009-5213-4098}.}
\email{daniel@halfounce.io}
\urladdr{https://halfounce.io}

\thanks{\emph{Computation and authorship.} All witnesses, checkers, and
enumerations in this note were produced by \textbf{Claude} (Anthropic),
directed by the author, on a single Apple~M4 laptop. The three shipped
checkers import only the Python standard library and share no code with
the search that produced the witnesses; the tournament and rigidity facts
were re-established here by code written from scratch. This addendum uses
no propositional solver: every claim is a finite check a reader can rerun.}

\thanks{\emph{Prior-art record.} This note answers Question~8.1 of the
companion determination \cite{PartG} of $k(3,4)=21$. The value was
absent from the literature before \cite{PartG}: the Erd\H{o}s problems
entry \cite{Blo} (\#112) records no exact $k(n,m)$, and \cite{IRW21}
remains the only paper on $r(I_m,L_n)$. A novelty sweep on the drafting
date, August~11, 2026, found nothing on the number of extremal
$\{I_3,\TT{4}\}$-free graphs on $20$ vertices, which is a fortiori new.
\emph{Draft of August~11, 2026.}}

\subjclass[2020]{Primary 05C55; Secondary 05C20, 05D10, 68V15}
\keywords{Oriented graph, Ramsey number, transitive tournament,
independent set, extremal graph, Paley tournament, rigidity, certified
computation}

\date{August 11, 2026}

\begin{abstract}
The companion note \cite{PartG} determines the Erd\H{o}s--Rado oriented
Ramsey number $k(3,4)=r(I_3,L_4)=21$ by an explicit $\{I_3,\TT{4}\}$-free
oriented graph $W$ on $20$ vertices together with a certified exhaustion
at $21$ vertices, and asks (Question~8.1) whether $W$ is the unique such
graph. We answer in the negative. There are at least thirteen pairwise
non-isomorphic $\{I_3,\TT{4}\}$-free oriented graphs on $20$ vertices,
and every one of the thirteen is \emph{rigid} (trivial automorphism
group), so the extremal configuration is far from a single symmetric
object. We record the local structure they share as a
short forcing lemma: in any $\{I_3,\TT{4}\}$-free oriented graph a vertex
has at most $7$ non-neighbours, and a vertex with exactly $7$ carries a
copy of the Paley tournament $\QR=\Cay(\mathbb{Z}_7,\{1,2,4\})$ as its
non-neighbourhood. Finally the algebraic route falls short: the largest
$\{I_3,\TT{4}\}$-free tournament blow-up has only $14$ vertices, six short
of the extremal order, and the thirteen exhibited witnesses are rigid
rather than Cayley. We close with the sharpened question of whether a
$\QR$ block is unavoidable in every witness.
\end{abstract}

\maketitle

\section{The question}\label{sec:q}

An \emph{oriented graph} is a loopless digraph with at most one arc
between any two vertices. Write $\I{n}$ for the independent set on $n$
vertices and $\TT{m}$ (equivalently $L_m$) for the transitive tournament
on $m$ vertices; call an oriented graph \emph{free} if it contains
neither $\I{3}$ nor $\TT{4}$. For a vertex $w$ let $N^{+}(w)$, $N^{-}(w)$,
$I(w)$ be its out-neighbours, in-neighbours, and non-neighbours, and call
$(p,q,s)=(|N^{+}(w)|,|N^{-}(w)|,|I(w)|)$ the \emph{type} of $w$, so
$p+q+s=|V|-1$. The companion note \cite{PartG} proves
$k(3,4)=r(I_3,L_4)=21$: every oriented graph on $21$ vertices contains an
$\I{3}$ or a $\TT{4}$, and the explicit $20$-vertex graph $W$ of its
Table~1 is free. Its Question~8.1 asks:

\begin{quote}
\emph{Is $W$ the unique free oriented graph on $20$ vertices up to
isomorphism? If not, what is the number of extremal graphs, and does any
admit an algebraic description?}
\end{quote}

The first part has a clean negative answer.

\begin{theorem}\label{thm:main}
There are at least thirteen pairwise non-isomorphic free oriented graphs
on $20$ vertices. Every one of the thirteen is rigid: its automorphism
group is trivial. In particular $W$ is not the unique $20$-vertex
extremal graph.
\end{theorem}

The thirteen graphs are $W$ itself together with twelve further free
graphs; all thirteen are printed as arc lists with the artifacts. Each of
the thirteen is in particular not vertex-transitive, since a
vertex-transitive graph on $n\ge 2$ vertices has $|\Aut|\ge n$. We
prove Theorem~\ref{thm:main} by direct verification
(\S\ref{sec:witnesses}); the local structure that the thirteen share is
explained by the forcing lemma of \S\ref{sec:qr7}; and the sense in which
the algebraic construction falls short is made precise in
\S\ref{sec:blowup}.

\section{The witnesses}\label{sec:witnesses}

Table~\ref{tab:inv} lists the thirteen graphs by three isomorphism
invariants: the number of arcs, the degree sequence of the
non-adjacency graph $\overline{N}$ (whose edges are the non-adjacent
pairs, so the $\overline{N}$-degree of $w$ is $s=|I(w)|$), and the number
of vertices of type $s=7$. By Lemma~\ref{lem:forcing} below the last
column equals the number of induced $\QR$ blocks.

\begin{table}[h]
\centering
\begin{tabular}{lccc}
\toprule
witness & arcs & $\overline{N}$-degree sequence & $\#\{s{=}7\}=\#\QR$ blocks\\
\midrule
$w_{1}=W$ & $126$ & $7^{11}\,6^{6}\,5^{3}$ & $11$\\
$w_{2}$   & $125$ & $7^{12}\,6^{6}\,5^{2}$ & $12$\\
$w_{3}$   & $128$ & $7^{6}\,6^{12}\,5^{2}$ & $6$\\
$w_{4}$   & $125$ & $7^{12}\,6^{6}\,5^{2}$ & $12$\\
$w_{5}$   & $126$ & $7^{11}\,6^{6}\,5^{3}$ & $11$\\
$w_{6}$   & $123$ & $7^{14}\,6^{6}$        & $14$\\
$w_{7}$   & $126$ & $7^{9}\,6^{10}\,5^{1}$ & $9$\\
$w_{8}$   & $126$ & $7^{12}\,6^{4}\,5^{4}$ & $12$\\
$w_{9}$   & $124$ & $7^{14}\,6^{4}\,5^{2}$ & $14$\\
$w_{10}$  & $127$ & $7^{8}\,6^{10}\,5^{2}$ & $8$\\
$w_{11}$  & $124$ & $7^{13}\,6^{6}\,5^{1}$ & $13$\\
$w_{12}$  & $125$ & $7^{13}\,6^{4}\,5^{3}$ & $13$\\
$w_{13}$  & $125$ & $7^{11}\,6^{8}\,5^{1}$ & $11$\\
\bottomrule
\end{tabular}
\caption{The thirteen free graphs on $20$ vertices. This coarse
fingerprint does not by itself separate all thirteen --- $w_1$ and $w_5$
agree in every column, as do $w_2$ and $w_4$ --- and the third column is
determined by the second. Separation is established by the canonical forms
and by the richer invariant the checker computes; see the certification. Arc counts range over
$123$--$128$ and the number of $\QR$ blocks over $6$--$14$.}
\label{tab:inv}
\end{table}

Each $w_i$ was verified free by an exhaustive test over all
$\binom{20}{3}=1140$ triples (none independent) and all
$\binom{20}{4}=4845$ quadruples (none a $\TT{4}$: a quadruple is a
$\TT{4}$ exactly when its six pairs are all adjacent and its four in-set
out-degrees are $0,1,2,3$). Rigidity, $|\Aut(w_i)|=1$, was certified two
independent ways: Weisfeiler--Leman colour refinement reaches twenty
distinct colours, so any automorphism fixes every vertex; and an explicit
automorphism backtracker returns count $1$. Pairwise non-isomorphism was
decided by a canonical form read off the discrete refinement, and
cross-checked by the richer invariant the checker computes; no two are
isomorphic.

\begin{remark}
The thirteen are only a lower bound on the extremal population. They were
obtained by completing random partial orientations to free graphs and
discarding isomorphic duplicates; in that sampling every completion that
survived deduplication was a new isomorphism class, which suggests --- but
does not prove --- that the family is large. Theorem~\ref{thm:main}
claims only what the thirteen exhibited certificates establish.
\end{remark}

\section{A forced Paley block}\label{sec:qr7}

The thirteen witnesses have no global symmetry, yet they share a rigid
local feature, visible in the last column of Table~\ref{tab:inv}: each
carries several induced copies of the Paley tournament. This is forced by
freeness alone.

\begin{lemma}[$\QR$ forcing]\label{lem:forcing}
Let $D$ be a free oriented graph and $w$ a vertex. Then $I(w)$ induces a
tournament with no transitive subtournament on $4$ vertices; consequently
$|I(w)|\le 7$, and if $|I(w)|=7$ then $I(w)$ induces a tournament
isomorphic to the Paley tournament $\QR=\Cay(\mathbb{Z}_7,\{1,2,4\})$.
\end{lemma}

\begin{proof}
If two vertices $x,y\in I(w)$ were non-adjacent, then $\{w,x,y\}$ would be
an $\I{3}$; so every pair inside $I(w)$ is adjacent and $I(w)$ induces a
tournament. A transitive subtournament on $4$ vertices inside $I(w)$ is a
$\TT{4}$ of $D$ (for a tournament target, induced and subdigraph
containment coincide), so $I(w)$ is $\TT{4}$-free.

It remains to bound and identify a $\TT{4}$-free tournament $T$. Every
tournament on $8$ vertices contains a transitive subtournament on $4$
vertices --- the tournament Ramsey value is $v(4)=8$ --- so $|V(T)|\le 7$;
and the unique tournament on $7$ vertices with no transitive
subtournament on $4$ vertices is the quadratic-residue tournament $\QR$
\cite{SanchezFlores98,ErdosMoser64}. Applying this to $T=I(w)$ gives
$|I(w)|\le 7$, with equality forcing $I(w)\cong\QR$.
\end{proof}

Both tournament inputs were re-established here from scratch, by
exhaustive enumeration: there are exactly $240$ labelled $\TT{4}$-free
tournaments on $7$ vertices, and every one is isomorphic to $\QR$ (so the
class is unique; and $240=7!/|\Aut(\QR)|$ with $|\Aut(\QR)|=21$), while
there are none on $8$ vertices. This matches the certified block
inventory of \cite[Prop.~3.3]{PartG} and the classical characterisation
\cite{SanchezFlores98}.

\begin{corollary}\label{cor:s7}
In any free oriented graph every vertex has at most $7$ non-neighbours;
and a vertex has exactly $7$ non-neighbours if and only if its
non-neighbourhood is a copy of $\QR$.
\end{corollary}

Corollary~\ref{cor:s7} is the endpoint $|I(w)|\le 7$ that drives the
counting bound of \cite[\S3]{PartG}. It also explains the last column of
Table~\ref{tab:inv}: in each witness the number of induced $\QR$ blocks
equals the number of type-$s{=}7$ vertices, which was confirmed
independently. Every one of the thirteen has between $6$ and $14$ such
vertices, so $\QR$ appears in each; whether it must appear in \emph{every}
free graph on $20$ vertices is left open in \S\ref{sec:open}.

\section{The algebraic route falls short}\label{sec:blowup}

The natural algebraic construction from $\QR$ is a blow-up, and it falls
well short of the extremal order.

\begin{proposition}\label{prop:blowup}
Let $T$ be a tournament and $m\ge 1$, and let $T[\I{m}]$ be the
lexicographic blow-up in which each vertex of $T$ is replaced by an
independent $m$-set. Then $T[\I{m}]$ is free if and only if $m\le 2$ and
$T$ is $\TT{4}$-free. Consequently the largest free graph of the form
$T[\I{m}]$ is $\QR[\I{2}]$ on $14$ vertices, which gives only
$k(3,4)\ge 15$.
\end{proposition}

\begin{proof}
In $T[\I{m}]$ two vertices are adjacent exactly when they lie in
different groups (within a group there are no arcs; between groups the
arc follows $T$, which is a tournament). Hence the underlying graph is
complete multipartite with parts of size $m$, its independence number is
$m$, and $T[\I{m}]$ is $\I{3}$-free iff $m\le 2$. A set of pairwise
adjacent vertices meets each group at most once, so it projects
injectively onto a subtournament of $T$ and the projection is an
isomorphism of induced sub-oriented-graphs; thus $T[\I{m}]$ contains a
$\TT{4}$ iff $T$ does. Combining the two, $T[\I{m}]$ is free iff $m\le 2$
and $T$ is $\TT{4}$-free. The largest $\TT{4}$-free tournament is $\QR$ on
$7$ vertices (Lemma~\ref{lem:forcing}), so the largest such blow-up is
$\QR[\I{2}]$ on $2\cdot 7=14$ vertices.
\end{proof}

The bound $k(3,4)\ge 15$ of Proposition~\ref{prop:blowup} is six short of
the truth $21$: the extremal order is $20$, but the best tournament
blow-up reaches $14$. The thirteen witnesses that reach the extremal
order are rigid (Theorem~\ref{thm:main}), hence not vertex-transitive,
hence not Cayley digraphs; and they are not blow-ups, since a proper
blow-up has two vertices with identical neighbourhoods and every $w_i$
has twenty distinct colour classes under refinement. In the language of
\cite[\S8.1]{PartG}, where the extremal graph for $k(3,3)=9$ is a single
circulant, the picture at $(3,4)$ is the opposite: the extremal graph is
not unique, and the thirteen exhibited here are aperiodic and
non-algebraic, with $\QR$ surviving only as a forced local block.

\section{Certification}\label{sec:cert}

The thirteen witnesses (as arc lists) and three standard-library checkers
are deposited with this note in the program's public certificate
repository \emph{Certify}
(\href{https://github.com/05oz/certify}{github.com/05oz/certify};
concept DOI \href{https://doi.org/10.5281/zenodo.21799111}{10.5281/zenodo.21799111}),
as Part~J. From each witness's arc list alone, \texttt{verify\_witnesses.py}
re-derives validity as an oriented graph, freeness over all
$\binom{20}{3}$ triples and $\binom{20}{4}$ quadruples, rigidity
(Weisfeiler--Leman refinement to twenty singleton colours, cross-checked
by an explicit automorphism count), and pairwise non-isomorphism (a
canonical form from the discrete refinement, cross-checked by the
invariant fingerprint of Table~\ref{tab:inv}); \texttt{verify\_qr7\_lemma.py}
enumerates the labelled $\TT{4}$-free tournaments on $7$ and $8$ vertices
($240$, all isomorphic to $\QR$; and $0$) and reports $|\Aut(\QR)|=21$;
and \texttt{blowup\_bound.py} checks $\QR[\I{2}]$ and $\QR[\I{3}]$. All
three checkers import only the Python standard library, share no code
with the search that produced the witnesses, and return \textsc{pass}.

\section{The sharpened question}\label{sec:open}

Question~8.1 of \cite{PartG} is answered: $W$ is not unique, and the
extremal count is at least $13$. The thirteen exhibited here are rigid,
hence none is vertex-transitive or Cayley; whether \emph{some other}
extremal graph might still be algebraic is not decided here. Three
questions remain.

\begin{question}\label{q:forced}
Is a $\QR$ block unavoidable? Equivalently, does every free oriented
graph on $20$ vertices have a vertex of type $s=7$?
\end{question}

By Corollary~\ref{cor:s7} a vertex of type $s=7$ is exactly a vertex
whose non-neighbourhood is a $\QR$; so Question~\ref{q:forced} asks
whether the propositional instance ``free on $20$ vertices with every
vertex having at most $6$ non-neighbours'' is unsatisfiable. All thirteen
exhibited witnesses carry between $6$ and $14$ such vertices, so $\QR$ is
present in each, but this is evidence, not a proof; the all-vertices-%
$s\le 6$ instance was not decided in a short trial and a certified answer
belongs to a dedicated computation.

\begin{question}
What is the exact number of free graphs on $20$ vertices? The thirteen
here are a lower bound; the sampling behaviour suggests a large family,
but no count is certified.
\end{question}

\begin{question}
Can the counting slack of \cite[\S8.1]{PartG} --- the cap of $24$ vertices
against the true maximum of $20$ --- be explained structurally, for
instance through the forced $\QR$ and Bermond blocks? Such an explanation
would bear on both the exact count and a human-readable upper bound.
\end{question}

\begin{thebibliography}{9}

\bibitem{Ber74} J.-C. Bermond, \emph{Some Ramsey numbers for directed
graphs}, Discrete Math. \textbf{9} (1974), 313--321.

\bibitem{Blo} T. F. Bloom, \emph{Erd\H{o}s Problem \#112},
\url{https://www.erdosproblems.com/112}, accessed 2026-08-11.

\bibitem{ErdosMoser64} P. Erd\H{o}s and L. Moser, \emph{On the
representation of directed graphs as unions of orderings}, Magyar Tud.
Akad. Mat. Kutat\'o Int. K\"ozl. \textbf{9} (1964), 125--132.

\bibitem{IRW21} F. Ihringer, D. Rajendraprasad and T. Weinert,
\emph{New bounds on the Ramsey number $r(I_m,L_n)$}, Discrete Math.
\textbf{344} (2021), no. 3, 112268. Also arXiv:1707.09556.

\bibitem{PartG} D. Kirtchakov, \emph{An Erd\H{o}s--Rado oriented Ramsey
number determined: $k(3,4)=r(I_3,L_4)=21$, by explicit witness and
certified exhaustion}, Certify (2026), version DOI
\href{https://doi.org/10.5281/zenodo.21890619}{10.5281/zenodo.21890619}.

\bibitem{SanchezFlores98} A. S\'anchez-Flores, \emph{On tournaments free
of large transitive subtournaments}, Graphs Combin. \textbf{14} (1998),
181--200.

\end{thebibliography}

\end{document}
